\chapter*{คำนำ}
\addcontentsline{toc}{chapter}{คำนำ}

เทคโนโลยีความจริงเสริม (Augmented Reality หรือ AR) และปัญญาประดิษฐ์ (Artificial Intelligence หรือ AI) ได้ก้าวขึ้นมาเป็นแกนกลางสำคัญในการขับเคลื่อนการปฏิวัติทางดิจิทัล ทั้งในภาคการศึกษา วิศวกรรม การแพทย์ และการสื่อสารเชิงพื้นที่ การผสานรวมการประมวลผลเชิงพื้นที่ (Spatial Computing) เข้ากับแบบจำลองวิทัศน์คอมพิวเตอร์และการเรียนรู้เชิงลึก (Deep Learning) เปิดโอกาสให้ผู้เรียนและนักพัฒนาสามารถสร้างสรรค์สภาพแวดล้อมการเรียนรู้เสมือนจริงที่เชื่อมต่อโลกกายภาพเข้ากับสารสนเทศดิจิทัลได้อย่างไร้รอยต่อ

คู่มือปฏิบัติการเรื่อง \textbf{การพัฒนาเทคโนโลยีความจริงเสริมร่วมกับปัญญาประดิษฐ์} เล่มนี้ เรียบเรียงขึ้นเพื่อใช้เป็นเอกสารประกอบการจัดการเรียนการสอนและการฝึกปฏิบัติการเชิงลึกสำหรับนักศึกษาระดับปริญญาตรีและผู้สนใจด้านการพัฒนาเทคโนโลยีเสมือนจริง โดยมุ่งเน้นการลงมือปฏิบัติจริง (Hands-on Active Learning) ผ่านเครื่องมือมาตรฐานเปิดระดับสากล ได้แก่ WebXR Device API, Three.js, A-Frame, MediaPipe Hands และ TensorFlow.js ซึ่งสามารถประมวลผลได้โดยตรงบนเว็บเบราว์เซอร์โดยไม่จำเป็นต้องติดตั้งซอฟต์แวร์ขนาดใหญ่

เนื้อหาภายในเล่มประกอบด้วย 8 หัวข้อปฏิบัติการหลัก เริ่มต้นตั้งแต่การวางโครงสร้างพื้นฐาน WebAR, การประมวลผลสัญญาณภาพและตรวจจับโครงกระดูกมือ 21 จุดร่วม, การคำนวณเวกเตอร์พิกัด 3 มิติและการตรวจจับท่าทางการจีบนิ้ว, การยิงลำแสงตรวจจับวัตถุเสมือน, การผนวกโมเดลการจำแนกภาพแบบ In-browser, การเชื่อมโยงสตรีมข้อมูล IoT สู่สภาพแวดล้อมเสมือน, การพัฒนาแบบจำลองห้องปฏิบัติการสะเต็มอัจฉริยะ, ไปจนถึงขั้นตอนการคอมไพล์เป็นแอปพลิเคชันมือถือเพื่อติดตั้งบนอุปกรณ์พกพาจริง

ผู้เขียนหวังเป็นอย่างยิ่งว่า คู่มือปฏิบัติการเล่มนี้จะมีส่วนช่วยส่งเสริมทักษะการคิดเชิงคำนวณ วิศวกรรมซอฟต์แวร์ และความคิดสร้างสรรค์แก่นักศึกษาและคณาจารย์ผู้สอน เพื่อร่วมกันยกระดับการศึกษาด้านเทคโนโลยีเกิดใหม่ของประเทศสู่ความเป็นเลิศอย่างยั่งยืน

\vspace{1.5cm}
\begin{flushright}
\textbf{ผศ.ดร.ชีวะ ทัศนา}\\
สาขาวิชาฟิสิกส์ คณะวิทยาศาสตร์และเทคโนโลยี\\
มหาวิทยาลัยราชภัฏรำไพพรรณี
\end{flushright}
